\chapter{Revisão bibliográfica}
\label{cap:revisao}

Este capítulo fundamenta as decisões que serão comparadas no experimento. A exposição distingue quatro níveis frequentemente tratados como equivalentes: HTTP é um protocolo de aplicação; REST é um estilo arquitetural; arquitetura orientada a eventos (EDA) organiza a colaboração entre componentes a partir da produção e reação a eventos; e Apache Kafka é uma plataforma que pode ser utilizada para implementar fluxos assíncronos \cite{fielding2002,aksakalli2021,kreps2011}. A distinção é necessária porque os efeitos observados não decorrem apenas da tecnologia, mas também das garantias adotadas pela aplicação, do momento em que uma operação é considerada concluída e do modo como falhas e repetições são tratadas. REST e EDA não são categorias mutuamente exclusivas: neste protótipo, ambas as variantes recebem requisições HTTP, e a diferença está no processamento imediato ou intermediado por mensageria.

\section{Registros e trilhas de auditoria}

Registros de auditoria descrevem ocorrências atribuíveis em um sistema e apoiam reconstrução cronológica, responsabilização e investigação posterior. Embora possam compartilhar infraestrutura com logs operacionais e de depuração, sua finalidade é diferente: um log de depuração auxilia o desenvolvimento, enquanto uma trilha de auditoria precisa conservar evidência coerente com a execução que pretende representar. \citeonline{amir2016} formalizam essa relação por propriedades de correção, completude e solidez do registro, mostrando que a análise posterior perde validade quando a instrumentação não corresponde adequadamente ao comportamento executado.

No protótipo, cada evento contém UUID, tipo do evento, tipo e identificador da entidade, ator, origem, instante e conteúdo contextual. O UUID permite correlacionar emissão, aceitação, consumo e persistência. Essa correlação também viabiliza verificar perda e duplicidade sem depender apenas dos totais agregados. O conteúdo é serializado de maneira determinística antes do cálculo de SHA-256. A representação denominada canônica neste trabalho ordena as propriedades do JSON e utiliza regras fixas de serialização no núcleo compartilhado; não se declara equivalência com um padrão universal de canonicalização. O mecanismo identifica conflitos entre o conteúdo recebido e o hash armazenado para um UUID. Não comprova a autoria do evento nem protege contra a alteração conjunta do registro e de seu hash por alguém com acesso ao banco.

Registros de auditoria são adequados ao experimento porque apresentam uma função pequena e repetível: receber, validar, identificar e persistir um fato. Essa simplicidade permite manter equivalência funcional entre as variantes. Em uso real, entretanto, a escolha arquitetural afeta o significado da confirmação. Uma resposta síncrona pode informar ao sistema de origem que o registro já está persistido. Em uma alternativa assíncrona, a aceitação pode ocorrer antes da persistência, exigindo mecanismos adicionais para acompanhar backlog, falhas e conclusão posterior.

\section{Comunicação em sistemas distribuídos}

Uma chamada entre processos não possui as mesmas garantias de uma chamada local. Rede, processo remoto e armazenamento podem falhar de modo independente, produzindo falhas parciais. O cliente pode receber timeout sem saber se a operação não iniciou, está em execução ou foi concluída antes de a resposta se perder. Por essa razão, repetição, timeout, identificação e observabilidade integram a semântica prática de uma operação distribuída.

Comunicação síncrona e assíncrona alteram onde essa incerteza aparece. Na interação síncrona, o cliente permanece temporalmente acoplado ao serviço e normalmente interpreta a resposta como conclusão do trabalho acordado. Na interação assíncrona, produtor e consumidor não precisam estar ativos no mesmo instante, mas a aplicação passa a distinguir aceitação, armazenamento intermediário e conclusão. O desacoplamento temporal reduz dependência simultânea, porém introduz estados intermediários que precisam ser operados e observados \cite{aksakalli2021,helland2007}.

\subsection{HTTP, REST e comunicação síncrona}

REST é um estilo arquitetural composto por restrições como cliente-servidor, ausência de estado de sessão no servidor, possibilidade de cache, interface uniforme e sistema em camadas \cite{fielding2002}. Portanto, uma API que utiliza HTTP e JSON não é automaticamente uma implementação completa de REST. \citeonline{pautasso2008} mostram que a escolha de serviços REST envolve decisões arquiteturais e consequências próprias, não apenas a seleção de um formato de mensagem. Neste estudo, a expressão ``variante REST'' identifica uma API orientada ao recurso de registro, utilizando HTTP de forma síncrona; a comparação não pretende avaliar individualmente todas as restrições do estilo.

Na restrição \textit{stateless}, cada requisição deve conter as informações necessárias para ser compreendida, sem depender de contexto de sessão mantido entre requisições no servidor. Segundo \citeonline{fielding2002}, essa restrição melhora a visibilidade das interações, favorece recuperação de falhas parciais e facilita a liberação de recursos entre requisições. A ausência de estado conversacional não significa ausência de dados persistidos: o servidor pode consultar bancos, mas não deve depender de uma sessão implícita para interpretar a requisição.

Os autores também apresentam limitações: repetir contexto aumenta o volume transferido, enquanto manter o estado da aplicação no cliente reduz o controle do servidor sobre a sequência de interação \cite{fielding2002}. Aplicado ao protótipo, isso significa que a ausência de sessão não elimina gargalos no processamento ou no banco. Uma API sem sessão ainda pode ficar indisponível se todas as requisições dependerem do mesmo armazenamento.

No fluxo estudado, o cliente envia \texttt{POST /audit}; a API valida o contrato, calcula o hash, persiste e somente então responde \texttt{201 Created}. Essa definição oferece uma pós-condição direta: ao receber sucesso, o cliente sabe que existe um registro persistido. Em contrapartida, validação, processamento e banco integram o caminho crítico. A latência do banco é incorporada à latência observada pelo cliente, conexões permanecem ocupadas e a indisponibilidade de uma dependência impede a conclusão. Timeouts e novas tentativas também exigem idempotência, pois a ausência de resposta não demonstra que o efeito não ocorreu.

\citeonline{pautasso2008} relacionam decisões de integração a consequências de desenvolvimento e manutenção. No recorte deste projeto, a variante síncrona possui menos componentes a operar, enquanto a assíncrona acrescenta broker e consumidor. Essa diferença permite discutir complexidade operacional, mas não demonstra, por si só, menor custo financeiro. O monitoramento deverá acompanhar requisições, duração, respostas HTTP, timeouts e uso de recursos. Escalar a API, sem ampliar a capacidade da dependência saturada, pode apenas deslocar ou intensificar a fila de espera.

\subsection{Mensageria e comunicação assíncrona}

Mensageria introduz produtor, intermediário e consumidor. O produtor registra uma mensagem no broker sem executar diretamente a lógica do consumidor. A retenção permite que produtor e consumidor operem em ritmos distintos e que o processamento seja retomado após indisponibilidade. Esse mecanismo pode absorver rajadas, mas não elimina o trabalho: quando a taxa de entrada supera a taxa de processamento, forma-se backlog, aumentando a latência de conclusão e a necessidade de capacidade posterior.

As expressões \textit{at-most-once}, \textit{at-least-once} e \textit{exactly-once} precisam ser associadas ao limite em que são garantidas. A primeira admite ausência de entrega; a segunda admite reentrega; a terceira exige explicitar qual efeito é realizado uma única vez. Kafka organiza o consumo por posições no log, enquanto o efeito produzido em um banco externo constitui outra fronteira de consistência \cite{kreps2011,helland2007}. Se o consumidor persistir no PostgreSQL e falhar antes de confirmar o offset, a mensagem poderá ser entregue novamente. O projeto busca entrega pelo menos uma vez e efeito idempotente para mensagens confirmadas no broker: a restrição única sobre \texttt{event\_id} impede uma segunda linha, e o hash diferencia reentrega legítima de conteúdo conflitante. A garantia ponta a ponta depende também de confirmação de publicação, retenção e retomada do consumidor. A implementação e as verificações funcionais desses mecanismos, incluindo o encaminhamento de falhas terminais à DLQ antes da confirmação de origem, são apresentadas na Seção~\ref{sec:validacao-funcional}; essas verificações não demonstram garantia universal sob quaisquer falhas.

Idempotência possui desafios práticos. O identificador precisa ser estável desde a origem; a verificação e o efeito devem estar protegidos contra concorrência; a política para o mesmo identificador com conteúdo diferente deve ser explícita; e a retenção da informação de deduplicação precisa cobrir o período de reentrega. Apenas consultar e depois inserir não é suficiente sob concorrência, pois dois consumidores podem observar ausência simultaneamente. Uma restrição única no banco torna a decisão atômica no ponto de persistência.

O uso de mensageria também muda manutenção e operação. O contrato do evento passa a ser uma dependência semântica entre produtor e consumidores. Alterações incompatíveis podem permanecer no tópico e falhar somente quando consumidas. Backlog, idade da mensagem, reentregas, registros não processáveis, estado dos grupos e diferença entre offsets produzidos e consumidos tornam-se sinais operacionais relevantes. Em compensação, a retenção permite recuperar consumidores e adicionar novos leitores sem bloquear o produtor.

\section{Arquitetura orientada a eventos}

Um evento representa um fato relevante que já ocorreu, enquanto um comando expressa uma solicitação para que algo seja feito. Em EDA, componentes publicam eventos e outros componentes reagem a eles; os padrões de comunicação discutidos por \citeonline{aksakalli2021} ajudam a distinguir essa colaboração de uma invocação direta. Neste trabalho, o evento representa uma ocorrência da aplicação cliente, e a ação do consumidor é registrar essa ocorrência. Kafka não define sozinho a arquitetura: contratos, responsabilidades, tratamento de falhas e evolução precisam ser organizados em torno dos eventos.

Entre as vantagens potenciais estão desacoplamento temporal, extensão por novos consumidores e isolamento de rajadas. Um produtor pode continuar aceitando eventos enquanto um consumidor está temporariamente indisponível, desde que o broker permaneça disponível e haja capacidade de retenção. Novos consumidores podem reutilizar o fluxo sem modificar a API produtora. Esses benefícios são condicionais: contratos continuam acoplando semanticamente os participantes, e o atraso entre produção e consumo cria consistência eventual.

As limitações aparecem principalmente na compreensão do estado global. Uma resposta de aceitação não representa a conclusão do processo; falhas podem ser descobertas depois que o cliente encerrou a interação; e a ordem entre eventos depende da estratégia de particionamento. \citeonline{laigner2021} observaram que consistência, evolução de esquemas e interleaving de fluxos transferem responsabilidades para a camada de aplicação. O diagnóstico requer identificadores de correlação e combinação de logs, métricas e rastreamento. Em entrevistas com 28 profissionais, \citeonline{niedermaier2019} identificaram observabilidade e monitoramento como desafios centrais de sistemas distribuídos e destacaram a necessidade de estratégia, responsabilidades e correlação das evidências operacionais.

Em termos práticos, EDA pode ser apropriada para integração desacoplada, propagação de mudanças, telemetria e trabalhos que toleram conclusão posterior. Ela é menos adequada quando o usuário precisa de confirmação imediata do efeito ou quando a complexidade operacional não se justifica. Custos incluem infraestrutura do broker, armazenamento e tráfego, instrumentação, operação de consumidores, gestão de contratos e capacitação da equipe. Escalar consumidores aumenta capacidade apenas até limites impostos por partições, banco de dados e trabalho por mensagem.

\section{Apache Kafka}

Kafka é uma plataforma de mensageria baseada em log distribuído. Produtores escrevem registros em tópicos, os tópicos são divididos em partições e consumidores controlam posições identificadas por offsets. A retenção independente da leitura permite que consumidores processem o fluxo em ritmos distintos e, dentro do período configurado, retomem posições anteriores \cite{kreps2011}.

O artigo de \citeonline{kreps2011} fundamenta o desenho original, mas não serve como comprovação das funcionalidades de versões posteriores. Seus resultados de desempenho pertencem à implementação e ao ambiente de 2011. Neste estudo, parâmetros do cliente, confirmação da publicação e comportamento sob falhas dependem da configuração efetivamente executada e deverão ser verificados no protótipo.

A ordenação é garantida dentro de uma partição, não de todo o tópico. Usar a mesma chave direciona eventos relacionados para a mesma partição quando o particionamento permanece compatível, preservando a ordem relativa desses eventos. Entretanto, aumentar o número de partições, alterar chaves ou combinar eventos de entidades distintas pode mudar a ordem global percebida. Há também um compromisso entre ordem e paralelismo: uma única partição oferece ordem total no tópico, mas limita o número de consumidores ativos no mesmo grupo.

Offsets representam a posição de consumo, e sua confirmação não é equivalente à persistência do efeito em um banco externo. Confirmar antes do efeito cria risco de perda lógica; confirmar depois permite reentrega. Uma transação que abranja Kafka e PostgreSQL não é fornecida automaticamente. Neste estudo, o consumidor persiste primeiro e confirma depois. A unicidade por UUID torna a reexecução segura quanto à duplicidade de linhas, embora continue necessário contabilizar reentregas e distinguir conflitos.

Outro ponto é a pressão de retorno. Kafka permite acumular mensagens quando consumidores ficam lentos, mas o backlog não pode crescer indefinidamente. Retenção, capacidade de disco, tempo máximo aceitável até a conclusão e velocidade de drenagem precisam ser monitorados. No experimento, a quantidade de mensagens pendentes e o tempo de recuperação após a retomada do consumidor representarão esse comportamento.

\section{Desempenho, resiliência e observabilidade}

Latência é um intervalo entre marcos definidos, não uma propriedade única. Em ambas as variantes, a latência HTTP corresponde ao envio até a resposta, e a latência de conclusão utiliza o envio até o mesmo marco registrado pelo repositório após o commit no PostgreSQL. A resposta síncrona sucede a transação; a assíncrona confirma o aceite no broker. Comparar o \texttt{202} diretamente com o \texttt{201} responderia apenas qual API devolve controle mais cedo, não qual fluxo conclui primeiro. O Capítulo~\ref{cap:metodo} especifica relógios, resolução e conciliação desses marcos.

Média pode ocultar caudas. Por isso, serão apresentados mediana e percentis p95 e p99, que indicam limites abaixo dos quais se encontram, respectivamente, 95\% e 99\% das observações. Esses percentis são sensíveis a aquecimento, pausas do ambiente, saturação de conexões, filas internas, coleta de lixo e contenção no banco. Serão calculados sobre amostras válidas de cada repetição, mantendo resultados individuais para mostrar dispersão.

Throughput é a quantidade de registros efetivamente concluídos por unidade de tempo. Na comparação principal, conclusão será persistência, evitando contar apenas respostas \texttt{202}. A taxa oferecida pelo gerador será registrada separadamente da taxa concluída. Se a entrada superar a capacidade, o throughput pode permanecer estável enquanto backlog e latência aumentam; por isso, as três medidas precisam ser interpretadas em conjunto.

Taxa de erro será calculada como operações sem o resultado esperado divididas pelas operações tentadas. Erros HTTP, falhas de publicação, falhas de consumo e conflitos de integridade serão separados, pois possuem significados distintos. CPU e memória ajudam a interpretar eficiência, mas dependem do limite de recursos e da instrumentação. Na alternativa assíncrona, backlog, reentregas e tempo de drenagem completam a análise.

Neste estudo, resiliência indica a capacidade do sistema de manter sua função ou retomá-la após uma falha prevista. Durante a indisponibilidade do consumidor, a API pode continuar aceitando mensagens, mas isso não significa que o fluxo completo continue concluindo registros. Serão observados disponibilidade da entrada, quantidade aceita, quantidade persistida, perda, duplicidade, backlog máximo e tempo necessário para retornar ao estado estável.

Observabilidade afeta a própria capacidade de validar o experimento. Cada evento terá UUID propagado nos logs e armazenado no banco. Logs registrarão marcos de recebimento, publicação, consumo, tentativa de persistência e confirmação. Essa correlação permite calcular a latência de conclusão e investigar valores extremos. Em sistemas distribuídos, a observabilidade auxilia a reconstrução de interações e o diagnóstico, mas exige decisões técnicas e organizacionais para lidar com a complexidade e a dinâmica dos componentes \cite{niedermaier2019}.

\section{Trabalhos relacionados}

\citeonline{cabane2024} compararam uma aplicação implementada como monólito e como arquitetura orientada a eventos, utilizando medidas de recursos, tempo de resposta, throughput e tráfego. No ambiente estudado, a alternativa monolítica utilizou menos recursos e respondeu mais rapidamente. Esse resultado não estabelece superioridade universal. Para este TCC, interessa o cuidado em comparar a mesma função com múltiplas métricas. O recorte proposto concentra-se em duas configurações de registro e acrescenta a distinção entre aceitação e conclusão. Como Kafka e consumidor também alteram a topologia e o uso de recursos, os resultados caracterizarão essas configurações concretas, sem isolar um efeito universal de REST ou EDA.

\citeonline{kazanavicius2023} avaliaram REST, RabbitMQ, Kafka, gRPC e GraphQL durante decomposição de monólitos. A amplitude permite discutir critérios de seleção, mas também distribui a análise entre cinco mecanismos e um contexto de migração. O estudo sustenta que a escolha depende dos requisitos e não apenas de desempenho bruto. O presente trabalho reduz o número de alternativas para aprofundar os efeitos temporais e de falha de duas variantes equivalentes, sem avaliar decomposição organizacional ou migração de um monólito.

\citeonline{blinowski2022} implementaram quatro versões de uma aplicação, combinando monólito ou microsserviços com Java ou C\#.NET, e avaliaram ambientes local e Azure. Entre os resultados, o monólito teve melhor desempenho em uma máquina, e acrescentar instâncias além de certo ponto prejudicou a aplicação. A análise explicita que ambiente e dimensionamento condicionam os benefícios da distribuição. Para transferir esses resultados ao presente problema, seria necessário considerar diferenças de tecnologia, carga e topologia. O TCC mantém linguagem, contrato, processamento e banco comuns às variantes, mas reconhece o broker e o consumidor como componentes adicionais da configuração assíncrona.

\citeonline{aksakalli2021} revisaram 239 publicações e selecionaram 38 estudos primários, identificando três abordagens de implantação, sete padrões de comunicação, oito desafios de implantação e seis desafios de comunicação. A revisão demonstra que implantação e comunicação são decisões igualmente relevantes e que padrões diferentes atendem contextos distintos. Por ser um estudo secundário, não fornece uma medição direta do comportamento de REST e Kafka sob a mesma carga; ele apoia a seleção dos conceitos e dos desafios que precisam ser observados experimentalmente.

\citeonline{laigner2021} combinaram literatura, inspeção de aplicações abertas e consulta a mais de 120 profissionais e pesquisadores para caracterizar gestão de dados em microsserviços. Foram relatados desafios de consistência eventual, evolução de esquemas, interleaving de eventos, replicação e validação de restrições na aplicação. O estudo explica por que idempotência e integridade não podem ser tratadas apenas como propriedades do broker. Entretanto, seu objetivo é caracterizar práticas de dados, e não comparar desempenho e recuperação de dois fluxos equivalentes.

As bases conceituais também apresentam limites de aplicação. \citeonline{helland2007} é um artigo de posição publicado no CIDR: discute padrões de coordenação e mensagens diante da ausência de transações globais, sem constituir um benchmark de REST e Kafka. \citeonline{amir2016} investigam a correção da instrumentação de auditoria e apresentam uma aplicação ao OpenMRS; o presente protótipo não reproduz sua prova formal. Já \citeonline{niedermaier2019} oferecem evidência qualitativa de entrevistas, útil para selecionar preocupações de observabilidade, mas insuficiente para estimar latência ou capacidade. A natureza de cada contribuição determina como ela é utilizada nesta pesquisa.

Os estudos convergem em três aspectos: distribuição adiciona custos; resultados dependem do contexto; e garantias de dados e operação não surgem automaticamente da tecnologia. Diferem quanto ao objeto, às alternativas e às métricas. No conjunto analisado, a oportunidade específica deste trabalho é comparar uma confirmação síncrona após persistência com uma aceitação assíncrona seguida de conclusão posterior, usando o mesmo contrato e armazenamento, e observar não apenas carga normal, mas indisponibilidade e drenagem de backlog. Essa lacuna é apresentada como recorte da busca e dos estudos selecionados, não como afirmação de inexistência absoluta de trabalhos.

\section{Síntese}

A variante síncrona escolhida responde após persistir o registro, incorporando as dependências necessárias ao efeito em seu caminho crítico. A ausência de sessão favorece certas propriedades de REST, mas não define sozinha o momento da confirmação. Na variante assíncrona, a retenção em Kafka permite separar produção e consumo e acomodar diferenças temporárias de ritmo, exigindo acompanhamento de backlog, reentrega, ordem e conclusão posterior. O experimento será estruturado para medir essas duas configurações, manter a função de registro equivalente e interpretar conjuntamente latência, recursos e garantias de processamento.
